\usemodule[basicexam]

\starttext

\startluacode
local BEM = thirddata.basicexam
local JN  = token.getexpansion("\\jobname")

-- 是否输出「通过」用例的详情（false 时只报失败）
local VERBOSE = true

---------------------------------------------------------------------------
-- 1) 显示宽度计算（中文按 2 计，ASCII 按 1 计）
---------------------------------------------------------------------------
local function disp_width(s)
  local w, i, n = 0, 1, #s
  while i <= n do
    local b = s:byte(i)
    if b < 0x80 then w = w + 1; i = i + 1
    elseif b < 0xE0 then w = w + 1; i = i + 2
    elseif b < 0xF0 then w = w + 2; i = i + 3
    elseif b < 0xF8 then w = w + 2; i = i + 4
    else w = w + 1; i = i + 1 end
  end
  return w
end

-- 按显示宽度右侧补空格
local function pad_disp(s, width)
  s = tostring(s or "")
  local d = width - disp_width(s)
  if d < 0 then d = 0 end
  return s .. string.rep(" ", d)
end

---------------------------------------------------------------------------
-- 2) 临时替换 structures.sections.current，控制章节号
---------------------------------------------------------------------------
local real_sections_current = structures.sections.current
local mock_chapter          = nil
structures.sections.current = function()
  if mock_chapter == nil then return nil end
  return { numberdata = { numbers = { nil, mock_chapter } } }
end

---------------------------------------------------------------------------
-- 3) 截获 context 输出
---------------------------------------------------------------------------
local real_context = context
local captured     = nil
local function capture(fn)
  captured = ""
  context = function(s) captured = tostring(s or "") end
  local ok, err = pcall(fn)
  context = real_context
  if not ok then error(err) end
  return captured
end

---------------------------------------------------------------------------
-- 4) 断言库
---------------------------------------------------------------------------
local T = { pass = 0, fail = 0, fails = {} }
local current_group = ""

local function group(name, intent)
  current_group = name
  texio.write_nl("")
  texio.write_nl("── " .. name .. " ──")
  if intent then
    for line in intent:gmatch("[^\n]+") do
      texio.write_nl("   " .. line)
    end
  end
end

local function eq(actual, expected, msg)
  local a = tostring(actual)
  local b = tostring(expected)
  local ok = (a == b)
  if ok then
    T.pass = T.pass + 1
  else
    T.fail = T.fail + 1
    T.fails[#T.fails + 1] = string.format("[%s] %s (expect=<%s> actual=<%s>)",
      current_group, msg or "?", b, a)
  end
  if VERBOSE or not ok then
    texio.write_nl(string.format("  %-4s| %s| expect=<%s>  actual=<%s>",
      ok and "ok" or "FAIL", pad_disp(msg or "?", 40), b, a))
  end
end

---------------------------------------------------------------------------
-- 5) 每个 case 前重建
---------------------------------------------------------------------------
local counter = 100
local function fresh()
  counter = counter + 1
  mock_chapter = counter
  BEM.check_mode = false
  BEM.data_new = {}
  BEM.init_tree()
  return counter
end

---------------------------------------------------------------------------
-- 开始
---------------------------------------------------------------------------
texio.write_nl("")
texio.write_nl("=== basicexam tests ===")

---------------------------------------------------------------------------
-- 基础路径
---------------------------------------------------------------------------
group("1. 单题单答", [[
验证最基本流程：push 一道题 → 加答案和分数 → pop。
预期：id/深度/分数/答案/is_leaf/order 全部正确。]] )
do
  local ch = fresh()
  BEM.push_level("question")
  eq(BEM.get_current_id(), string.format("%s-%d-Q-1", JN, ch), "id")
  eq(BEM.get_depth(), 1, "depth")
  BEM.add_scored_answer("A", 5)
  BEM.pop_level("question")
  eq(BEM.get_depth(), 0, "回到根")
  eq(BEM.answer_tree.totalnumber, 1, "root totalnumber")
  eq(BEM.answer_tree.totalpoint, 5, "root totalpoint")
  local d = BEM.get_display(string.format("%s-%d-Q-1", JN, ch))
  eq(d.point, 5, "point")
  eq(d.answer, "A", "answer")
  eq(d.is_leaf, true, "is_leaf")
  eq(d.order, "1", "order")
end

group("2. 嵌套子题", [[
验证父子两层：父节点有子题后不再是叶子，
父分数 = 所有子分数之和，order 形成层级路径。]])
do
  local ch = fresh()
  BEM.push_level("q")
    BEM.push_level("s")
      BEM.add_scored_answer("x", 2)
    BEM.pop_level("s")
  BEM.pop_level("q")
  local Q = string.format("%s-%d-Q-1", JN, ch)
  local S = Q .. "-1"
  local q = BEM.get_display(Q)
  eq(q.is_leaf, false, "q is_leaf")
  eq(q.point, 2, "q totalpoint")
  eq(q.totalnumber, 1, "q totalnumber")
  eq(q.order, "1", "q order")
  local s = BEM.get_display(S)
  eq(s.is_leaf, true, "s is_leaf")
  eq(s.point, 2, "s point")
  eq(s.answer, "x", "s answer")
  eq(s.order, "1-1", "s order")
end

group("3. 多章节", [[
两章各一道题，验证章节号进入 id 且互不干扰。]])
do
  local ch1 = fresh()
  BEM.push_level("q"); BEM.add_scored_answer("a", 5); BEM.pop_level("q")
  mock_chapter = ch1 + 1000
  BEM.push_level("q"); BEM.add_scored_answer("b", 3); BEM.pop_level("q")
  eq(BEM.get_display(string.format("%s-%d-Q-1", JN, ch1)).point, 5, "ch1 point")
  eq(BEM.get_display(string.format("%s-%d-Q-1", JN, ch1 + 1000)).point, 3, "ch2 point")
end

group("4. 无章节", [[
用户完全不写 chapter 时，章节号应为 0，一切照常工作。]])
do
  mock_chapter = nil
  BEM.check_mode = false
  BEM.init_tree()
  BEM.push_level("q")
  eq(BEM.get_current_id(), string.format("%s-0-Q-1", JN), "无章节 id")
  BEM.add_scored_answer("a", 4)
  BEM.pop_level("q")
  eq(BEM.get_display(string.format("%s-0-Q-1", JN)).point, 4, "无章节 point")
end

group("5. 叶子转非叶子", [[
先给父节点加答案（成为叶子），再加子题。
此时父节点旧的答案和分数应被撤销，只保留子题之和。]])
do
  local ch = fresh()
  BEM.push_level("q")
  BEM.add_scored_answer("A", 5)
  BEM.push_level("s")
  BEM.add_scored_answer("B", 3)
  BEM.pop_level("s")
  BEM.pop_level("q")
  local Q = string.format("%s-%d-Q-1", JN, ch)
  eq(BEM.answer_tree.totalnumber, 1, "root totalnumber")
  eq(BEM.answer_tree.totalpoint, 3, "root totalpoint（旧 5 被替换）")
  eq(BEM.get_display(Q).point, 3, "q totalpoint")
  eq(BEM.get_display(Q).is_leaf, false, "q is_leaf")
  eq(BEM.get_display(Q .. "-1").point, 3, "sub point")
end

group("6. 深层 allleafnumber", [[
三层嵌套，其中一层有两个子题。
leafnumber 只统计「顶层叶子题」，allleafnumber 统计所有叶子。]])
do
  local ch = fresh()
  BEM.push_level("q")
    BEM.push_level("s")
      BEM.push_level("ss"); BEM.add_scored_answer("a", 1); BEM.pop_level("ss")
      BEM.push_level("ss"); BEM.add_scored_answer("b", 1); BEM.pop_level("ss")
    BEM.pop_level("s")
    BEM.push_level("s"); BEM.add_scored_answer("c", 1); BEM.pop_level("s")
  BEM.pop_level("q")
  BEM.save_structure_data()
  local s = BEM.get_structure_stats(1, ch)
  eq(s.totalnumber, 1, "顶层题数")
  eq(s.totalpoint, 3, "总分")
  eq(s.leafnumber, 0, "顶层非叶子数")
  eq(s.allleafnumber, 3, "叶子总数")
end

group("7. extra_data", [[
节点可携带额外元数据（如选择题选项数）。
验证存/取、缺失键、数字默认值三种情况。]])
do
  local ch = fresh()
  BEM.push_level("q")
  BEM.set_extra_data("totalitemnumber", "4")
  BEM.add_scored_answer("a", 1)
  BEM.pop_level("q")
  local id = string.format("%s-%d-Q-1", JN, ch)
  BEM.build_node_data()
  -- tex_get_data 读的是上一次编译的数据文件；
  -- 这里把本次生成的记录当作“下一次编译读到的数据”。
  local saved_data = BEM.data
  BEM.data = BEM.data_new
  eq(capture(function() BEM.tex_get_data(id, "totalitemnumber") end), "4", "extra data 读回")
  eq(capture(function() BEM.tex_get_data(id, "not_exist") end), "", "缺失键输出空串")
  eq(capture(function() BEM.tex_get_extra_number(id, "totalitemnumber", 0) end), "4", "extra number")
  eq(capture(function() BEM.tex_get_extra_number(id, "not_exist", 99) end), "99", "缺失时默认值")
  BEM.data = saved_data
end

---------------------------------------------------------------------------
-- 环境栈
---------------------------------------------------------------------------
group("8. 跨层闭合检测", [[
用 pop_level("q") 去关 "s" 环境时应被拒绝，栈保持不变。]])
do
  local ch = fresh()
  BEM.enable_check()
  BEM.push_level("q")
  BEM.push_level("s")
  local before = BEM.get_current_id()
  BEM.pop_level("q")
  eq(BEM.get_current_id(), before, "错误 pop 后栈未变")
  BEM.pop_level("s")
  BEM.pop_level("q")
  BEM.check_mode = false
end

group("9. 未闭合检测", [[
文档末尾扫描未闭合的环境，验证栈计数正确。]])
do
  local ch = fresh()
  BEM.push_level("q")
  BEM.push_level("s")
  eq(#BEM.level_stack, 3, "栈深 root+q+s")
  BEM.check_unclosed()
  BEM.pop_level("s")
  BEM.pop_level("q")
  eq(#BEM.level_stack, 1, "清空回到 root")
end

---------------------------------------------------------------------------
-- order / 输出
---------------------------------------------------------------------------
group("10. order 生成", [[
层级路径式 order：顶层=1，第二层=1-1 / 1-2，第三层=1-1-1 / 1-1-2。]])
do
  local ch = fresh()
  BEM.push_level("q")
    BEM.push_level("s")
      BEM.push_level("ss"); BEM.add_scored_answer("a", 1); BEM.pop_level("ss")
      BEM.push_level("ss"); BEM.add_scored_answer("b", 1); BEM.pop_level("ss")
    BEM.pop_level("s")
    BEM.push_level("s"); BEM.add_scored_answer("c", 1); BEM.pop_level("s")
  BEM.pop_level("q")
  local P = string.format("%s-%d-Q-1", JN, ch)
  eq(BEM.get_display(P).order, "1", "Q order")
  eq(BEM.get_display(P .. "-1").order, "1-1", "s1 order")
  eq(BEM.get_display(P .. "-1-1").order, "1-1-1", "ss1 order")
  eq(BEM.get_display(P .. "-1-2").order, "1-1-2", "ss2 order")
  eq(BEM.get_display(P .. "-2").order, "1-2", "s2 order")
end

group("11. tex_get 输出格式", [[
TeX 接口对不同 key 的输出格式：
point 整数不带小数点，小数保留 1 位；is_leaf 为 "true"/"false"。]])
do
  local ch = fresh()
  BEM.push_level("q"); BEM.add_scored_answer("A", 7.5); BEM.pop_level("q")
  BEM.push_level("q"); BEM.add_scored_answer("B", 5); BEM.pop_level("q")
  local id1 = string.format("%s-%d-Q-1", JN, ch)
  local id2 = string.format("%s-%d-Q-2", JN, ch)
  eq(capture(function() BEM.tex_get(id1, "point") end), "7.5", "小数 point")
  eq(capture(function() BEM.tex_get(id2, "point") end), "5", "整数 point")
  eq(capture(function() BEM.tex_get(id1, "answer") end), "A", "answer")
  eq(capture(function() BEM.tex_get(id1, "is_leaf") end), "true", "is_leaf")
  eq(capture(function() BEM.tex_get(id1, "order") end), "1", "order")
end

---------------------------------------------------------------------------
-- 统计
---------------------------------------------------------------------------
group("12. get_question_level_stats", [[
针对单个题目的层级统计：level1 子题数、叶子数、总分。]])
do
  local ch = fresh()
  BEM.push_level("q")
    BEM.push_level("s"); BEM.add_scored_answer("a", 1); BEM.pop_level("s")
    BEM.push_level("s"); BEM.add_scored_answer("b", 2); BEM.pop_level("s")
  BEM.pop_level("q")
  local st = BEM.get_question_level_stats(string.format("%s-%d-Q-1", JN, ch))
  eq(st.levelcounts[1], 2, "level1 子题数")
  eq(st.leafnumber, 2, "叶子数")
  eq(st.totalpoint, 3, "总分")
end

group("13. 未知 id 容错", [[
查询不存在的 id 应返回默认值而非报错。]])
do
  fresh()
  local d = BEM.get_display("nonexistent")
  eq(d.point, 0, "point")
  eq(d.answer, "", "answer")
  eq(d.is_leaf, false, "is_leaf")
end

group("14. save_structure_data 往返", [[
先 save 进 dataset，再读回，验证数据一致性。]])
do
  local ch = fresh()
  BEM.push_level("q")
    BEM.push_level("s"); BEM.add_scored_answer("a", 2); BEM.pop_level("s")
    BEM.push_level("s"); BEM.add_scored_answer("b", 3); BEM.pop_level("s")
  BEM.pop_level("q")
  BEM.save_structure_data()
  local s = BEM.get_structure_stats(1, ch)
  eq(s.totalnumber, 1, "totalnumber")
  eq(s.totalpoint, 5, "totalpoint")
  eq(s.leafnumber, 0, "leafnumber")
  eq(s.allleafnumber, 2, "allleafnumber")
end

group("15. tex_get_structure / tex_get_question_level", [[
TeX 层接口：按结构层和按题目 id 输出统计。]])
do
  local ch = fresh()
  BEM.push_level("q")
    BEM.push_level("s"); BEM.add_scored_answer("a", 2); BEM.pop_level("s")
    BEM.push_level("s"); BEM.add_scored_answer("b", 3); BEM.pop_level("s")
  BEM.pop_level("q")
  BEM.save_structure_data()
  eq(capture(function() BEM.tex_get_structure(1, ch, "totalnumber")   end), "1", "str totalnumber")
  eq(capture(function() BEM.tex_get_structure(1, ch, "totalpoint")    end), "5", "str totalpoint")
  eq(capture(function() BEM.tex_get_structure(1, ch, "allleafnumber") end), "2", "str allleafnumber")
  local id = string.format("%s-%d-Q-1", JN, ch)
  eq(capture(function() BEM.tex_get_question_level(id, "level1")     end), "2", "level1")
  eq(capture(function() BEM.tex_get_question_level(id, "leafnumber") end), "2", "leafnumber")
  eq(capture(function() BEM.tex_get_question_level(id, "totalpoint") end), "5", "totalpoint")
end

group("16. 章节切换 order 隔离", [[
两章各一道题，order 都应为 1（按章节重置）。]])
do
  mock_chapter = 400
  BEM.check_mode = false
  BEM.init_tree()
  BEM.push_level("q"); BEM.pop_level("q")
  mock_chapter = 401
  BEM.push_level("q"); BEM.pop_level("q")
  eq(BEM.get_display(string.format("%s-400-Q-1", JN)).order, "1", "ch400 order")
  eq(BEM.get_display(string.format("%s-401-Q-1", JN)).order, "1", "ch401 order")
end

---------------------------------------------------------------------------
-- 分支覆盖
---------------------------------------------------------------------------
group("17. get_structure_stats 无缓存路径", [[
关键分支：不先 save，直接读统计时走「遍历 answer_tree 重算」路径。
这是首次编译时真正的执行路径。]])
do
  local ch = fresh()
  BEM.push_level("q")
    BEM.push_level("s"); BEM.add_scored_answer("a", 2); BEM.pop_level("s")
    BEM.push_level("s"); BEM.add_scored_answer("b", 3); BEM.pop_level("s")
  BEM.pop_level("q")
  local s = BEM.get_structure_stats(1, ch)
  eq(s.totalnumber, 1, "无缓存 totalnumber")
  eq(s.totalpoint, 5, "无缓存 totalpoint")
  eq(s.leafnumber, 0, "无缓存 leafnumber")
  eq(s.allleafnumber, 2, "无缓存 allleafnumber")
end

group("18. 父节点 answer 拼接", [[
父节点非叶子时，answer 应为所有子答案用 "; " 拼接。]])
do
  local ch = fresh()
  BEM.push_level("q")
    BEM.push_level("s"); BEM.add_scored_answer("a", 1); BEM.pop_level("s")
    BEM.push_level("s"); BEM.add_scored_answer("b", 1); BEM.pop_level("s")
    BEM.push_level("s"); BEM.add_scored_answer("c", 1); BEM.pop_level("s")
  BEM.pop_level("q")
  local q = BEM.get_display(string.format("%s-%d-Q-1", JN, ch))
  eq(q.answer, "a; b; c", "三个子答案拼接")
  eq(q.is_leaf, false, "q 不是叶子")
  eq(q.point, 3, "q 总分")
end

group("19. 重复 add_scored_answer", [[
同一节点二次给分应「覆盖」而非「累加」：
旧分先从祖先撤销，再写入新分。]])
do
  local ch = fresh()
  BEM.push_level("q")
  BEM.add_scored_answer("A", 5)
  BEM.add_scored_answer("B", 3)
  BEM.pop_level("q")
  local d = BEM.get_display(string.format("%s-%d-Q-1", JN, ch))
  eq(d.point, 3, "point 覆盖为 3")
  eq(d.answer, "B", "answer 覆盖为 B")
  eq(BEM.answer_tree.totalpoint, 3, "root 不重复累加")
  eq(BEM.answer_tree.totalnumber, 1, "题目数仍为 1")
end

group("20. 同章 Q-N 递增", [[
同章节三道题的 id 分别为 Q-1/Q-2/Q-3，root 计数累加。]])
do
  local ch = fresh()
  BEM.push_level("q"); BEM.add_scored_answer("a", 1); BEM.pop_level("q")
  BEM.push_level("q"); BEM.add_scored_answer("b", 2); BEM.pop_level("q")
  BEM.push_level("q"); BEM.add_scored_answer("c", 3); BEM.pop_level("q")
  local id1 = string.format("%s-%d-Q-1", JN, ch)
  local id2 = string.format("%s-%d-Q-2", JN, ch)
  local id3 = string.format("%s-%d-Q-3", JN, ch)
  eq(BEM.get_display(id1).point, 1, "Q-1 分数")
  eq(BEM.get_display(id2).point, 2, "Q-2 分数")
  eq(BEM.get_display(id3).point, 3, "Q-3 分数")
  eq(BEM.answer_tree.totalnumber, 3, "root totalnumber")
  eq(BEM.answer_tree.totalpoint, 6, "root totalpoint")
end

group("21. 多余 pop 边界", [[
已在 root 时再 pop，应静默返回、不崩、栈不变。]])
do
  fresh()
  eq(#BEM.level_stack, 1, "初始只有 root")
  BEM.pop_level()
  eq(#BEM.level_stack, 1, "多余 pop 后栈不变")
  eq(BEM.get_current_id(), "root", "仍在 root")
  eq(BEM.get_depth(), 0, "depth 仍为 0")
end

group("22. 根节点拒绝加答案", [[
在 root（depth=0）上 add_scored_answer 应被拒绝。]])
do
  fresh()
  eq(BEM.get_depth(), 0, "当前在 root")
  BEM.add_scored_answer("X", 99)
  eq(BEM.answer_tree.totalpoint, 0, "root 分数仍为 0")
  eq(BEM.answer_tree.totalnumber, 0, "root 题数仍为 0")
  eq(BEM.answer_tree.is_leaf, false, "root 不是叶子")
  eq(BEM.answer_tree.answer, "", "root answer 空")
end

group("23. 有子节点时拒绝加答案", [[
已有子题的父节点不能再加自己的答案（防止污染子题之和）。]])
do
  local ch = fresh()
  BEM.push_level("q")
    BEM.push_level("s"); BEM.add_scored_answer("child", 2); BEM.pop_level("s")
  BEM.add_scored_answer("should_be_ignored", 99)
  BEM.pop_level("q")
  local q = BEM.get_display(string.format("%s-%d-Q-1", JN, ch))
  eq(q.point, 2, "q 总分仍为子题之和")
  eq(q.is_leaf, false, "q 仍不是叶子")
  eq(BEM.answer_tree.totalpoint, 2, "root 总分未受污染")
end

---------------------------------------------------------------------------
-- 极端情况
---------------------------------------------------------------------------
group("24. 浮点累加精度", [[
0.1 + 0.2 在二进制浮点下并不等于 0.3。
但 tex_get 用 %.1f 格式化后应输出 "0.3"。]])
do
  local ch = fresh()
  BEM.push_level("q")
    BEM.push_level("s"); BEM.add_scored_answer("a", 0.1); BEM.pop_level("s")
    BEM.push_level("s"); BEM.add_scored_answer("b", 0.2); BEM.pop_level("s")
  BEM.pop_level("q")
  local id = string.format("%s-%d-Q-1", JN, ch)
  eq(capture(function() BEM.tex_get(id, "point") end), "0.3", "格式化后 0.3")
end

group("25. 非法输入容错", [[
空答案 / nil 答案 / nil 分 / 非数字字符串分 / 字符串数字分。
预期：空答案保留空，非法分退化为 0，字符串数字能正确解析。]])
do
  local ch = fresh()
  BEM.push_level("q"); BEM.add_scored_answer("", 5); BEM.pop_level("q")
  eq(BEM.get_display(string.format("%s-%d-Q-1", JN, ch)).point, 5, "空答案保留 5 分")

  ch = fresh()
  BEM.push_level("q"); BEM.add_scored_answer(nil, 5); BEM.pop_level("q")
  eq(BEM.get_display(string.format("%s-%d-Q-1", JN, ch)).answer, "", "nil 答案存空串")

  ch = fresh()
  BEM.push_level("q"); BEM.add_scored_answer("a", nil); BEM.pop_level("q")
  eq(BEM.get_display(string.format("%s-%d-Q-1", JN, ch)).point, 0, "nil 分 → 0")

  ch = fresh()
  BEM.push_level("q"); BEM.add_scored_answer("a", "abc"); BEM.pop_level("q")
  eq(BEM.get_display(string.format("%s-%d-Q-1", JN, ch)).point, 0, "非数字分 → 0")

  ch = fresh()
  BEM.push_level("q"); BEM.add_scored_answer("a", "5"); BEM.pop_level("q")
  eq(BEM.get_display(string.format("%s-%d-Q-1", JN, ch)).point, 5, "字符串数字 5 解析")
end

group("26. 答案含分隔符和特殊字符", [[
子答案内部含逗号时不影响拼接；分隔符是 "; "。
当前实现不转义，按原样拼接。]])
do
  local ch = fresh()
  BEM.push_level("q")
    BEM.push_level("s"); BEM.add_scored_answer("a,b", 1); BEM.pop_level("s")
    BEM.push_level("s"); BEM.add_scored_answer("c d", 1); BEM.pop_level("s")
  BEM.pop_level("q")
  local q = BEM.get_display(string.format("%s-%d-Q-1", JN, ch))
  eq(q.answer, "a,b; c d", "拼接按原样不转义")
end

group("27. 极深嵌套 10 层", [[
十层 push → 内层加一次答案 → 逐层 pop。
验证递归/栈深度不会失控，root 计数正确。]])
do
  local ch = fresh()
  for _ = 1, 10 do BEM.push_level("q") end
  eq(#BEM.level_stack, 11, "栈深 root+10")
  BEM.add_scored_answer("deep", 1)
  for _ = 1, 10 do BEM.pop_level("q") end
  eq(#BEM.level_stack, 1, "回到 root")
  eq(BEM.answer_tree.totalnumber, 1, "root totalnumber")
  eq(BEM.answer_tree.totalpoint, 1, "root totalpoint")
end

group("28. section 结构残缺", [[
sectiondata.numberdata.numbers 为空表时，
get_current_chapter 应容错为 0（无章节）。]])
do
  -- 直接改 mock，返回带空 numbers 的假 section
  local saved = structures.sections.current
  structures.sections.current = function()
    return { numberdata = { numbers = {} } }
  end
  BEM.check_mode = false
  BEM.init_tree()
  BEM.push_level("q")
  eq(BEM.get_current_id(), string.format("%s-0-Q-1", JN), "空 numbers → 章节 0")
  BEM.add_scored_answer("a", 1)
  BEM.pop_level("q")
  structures.sections.current = saved
end

group("29. 章节号为字符串", [[
mock 章节号为 "5"（字符串），
get_current_chapter 应用 tonumber 正确转换。]])
do
  local saved_mock = mock_chapter
  mock_chapter = "5"
  BEM.check_mode = false
  BEM.init_tree()
  BEM.push_level("q"); BEM.add_scored_answer("a", 1); BEM.pop_level("q")
  eq(BEM.get_display(string.format("%s-5-Q-1", JN)).point, 1, "字符串章节号可用")
  mock_chapter = saved_mock
end

group("30. push_level 无 env_type", [[
env_type 为 nil 时应存为 "unknown"，不崩溃。]])
do
  fresh()
  BEM.push_level()
  eq(BEM.current_node.env_type, "unknown", "默认 env_type")
  BEM.add_scored_answer("a", 1)
  BEM.pop_level()
  eq(BEM.answer_tree.totalpoint, 1, "root 分数")
end

group("31. 当前 structure 统计", [[
tex_get_structure 不带 number 时走「当前 section」路径，
等价于 UnlimitedCurrentStructureTotalNumber。]])
do
  local ch = fresh()   -- mock_chapter = ch
  BEM.push_level("q")
    BEM.push_level("s"); BEM.add_scored_answer("a", 3); BEM.pop_level("s")
  BEM.pop_level("q")
  BEM.save_structure_data()
  eq(capture(function() BEM.tex_get_structure(1, "", "totalnumber")   end), "1", "当前 totalnumber")
  eq(capture(function() BEM.tex_get_structure(1, "", "totalpoint")    end), "3", "当前 totalpoint")
  eq(capture(function() BEM.tex_get_structure(1, "", "allleafnumber") end), "1", "当前 allleafnumber")
end

group("32. get_current_chapter 缓存命中", [[
同一 sectiondata 被返回两次时，应走缓存路径并给出相同章节号。
（真实 ConTeXt 在同一章内就是这样。）]])
do
  local shared = { numberdata = { numbers = { nil, 777 } } }
  local saved = structures.sections.current
  structures.sections.current = function() return shared end
  BEM.check_mode = false
  BEM.init_tree()
  BEM.push_level("q"); BEM.pop_level("q")   -- 首次：缓存 miss，记 777
  BEM.push_level("q"); BEM.pop_level("q")   -- 再次：缓存 hit
  eq(BEM.get_display(string.format("%s-777-Q-1", JN)) ~= nil, true, "第一次 Q-1 存在")
  eq(BEM.get_display(string.format("%s-777-Q-2", JN)) ~= nil, true, "第二次 Q-2 存在")
  structures.sections.current = saved
end

group("33. init_tree 清空 node_index", [[
重建后，上一次的 id 不应还能被 get_display 查到。]])
do
  local ch = fresh()
  BEM.push_level("q"); BEM.add_scored_answer("a", 1); BEM.pop_level("q")
  local old_id = string.format("%s-%d-Q-1", JN, ch)
  eq(BEM.get_display(old_id).point, 1, "重建前可查")
  BEM.init_tree()
  eq(BEM.get_display(old_id).point, 0, "重建后查不到（返回默认 0）")
end

group("34. 同章 order 递增", [[
同章三道题的 order 应为 1、2、3。]])
do
  local ch = fresh()
  for _ = 1, 3 do
    BEM.push_level("q"); BEM.add_scored_answer("x", 1); BEM.pop_level("q")
  end
  local P = string.format("%s-%d-Q-", JN, ch)
  eq(BEM.get_display(P .. "1").order, "1", "Q-1 order")
  eq(BEM.get_display(P .. "2").order, "2", "Q-2 order")
  eq(BEM.get_display(P .. "3").order, "3", "Q-3 order")
end

group("35. 负分处理", [[
负分是否允许？当前实现会接受并累加，确认这是你要的语义。]])
do
  local ch = fresh()
  BEM.push_level("q"); BEM.add_scored_answer("penalty", -2); BEM.pop_level("q")
  local d = BEM.get_display(string.format("%s-%d-Q-1", JN, ch))
  eq(d.point, -2, "叶子分数为 -2")
  eq(BEM.answer_tree.totalpoint, -2, "root 总分 -2")
end

group("36. level ≥ 2 明确拒绝", [[
查询 level=3 时，模块应明确返回 0（不支持），
而不是像旧版那样静默回退到章节统计冒充。]])
do
  local saved = structures.sections.current
  structures.sections.current = function()
    return { numberdata = { numbers = { nil, 555 } } }
  end
  BEM.check_mode = false
  BEM.init_tree()
  BEM.push_level("q"); BEM.add_scored_answer("a", 1); BEM.pop_level("q")
  BEM.save_structure_data()
  local s = BEM.get_current_structure_stats(3)
  eq(s.totalnumber, 0, "level 3 返回 0（明确拒绝）")
  structures.sections.current = saved
end



group("37. jobname 含连字符", [[
jobname = "my-doc" 时 id 变成 my-doc-101-Q-1，
解析和查找都不受影响。]])
do
  local saved_expansion = token.getexpansion
  local saved_current   = structures.sections.current
  token.getexpansion = function(name)
    if name == "\\jobname" then return "my-doc" end
    return saved_expansion(name)
  end
  structures.sections.current = function()
    return { numberdata = { numbers = { nil, 101 } } }
  end

  BEM.check_mode = false
  BEM.init_tree()
  BEM.push_level("q")
  BEM.add_scored_answer("a", 1)
  BEM.pop_level("q")

  eq(BEM.get_display("my-doc-101-Q-1").point, 1, "可查")
  eq(BEM.get_display("my-doc-101-Q-1").order, "1", "order 正常")

  structures.sections.current = saved_current
  token.getexpansion = saved_expansion
end
---------------------------------------------------------------------------
-- 收尾
---------------------------------------------------------------------------
structures.sections.current = real_sections_current
BEM.check_mode = false
BEM.init_tree()

texio.write_nl("")
texio.write_nl(string.format("=== %d passed, %d failed ===", T.pass, T.fail))
if T.fail > 0 then
  for _, f in ipairs(T.fails) do
    texio.write_nl("  - " .. f)
  end
end
\stopluacode

\stoptext